% Options for packages loaded elsewhere
\PassOptionsToPackage{unicode}{hyperref}
\PassOptionsToPackage{hyphens}{url}
%
\documentclass[
]{article}
\usepackage{amsmath,amssymb}
\usepackage{iftex}
\ifPDFTeX
  \usepackage[T1]{fontenc}
  \usepackage[utf8]{inputenc}
  \usepackage{textcomp} % provide euro and other symbols
\else % if luatex or xetex
  \usepackage{unicode-math} % this also loads fontspec
  \defaultfontfeatures{Scale=MatchLowercase}
  \defaultfontfeatures[\rmfamily]{Ligatures=TeX,Scale=1}
\fi
\usepackage{lmodern}
\ifPDFTeX\else
  % xetex/luatex font selection
\fi
% Use upquote if available, for straight quotes in verbatim environments
\IfFileExists{upquote.sty}{\usepackage{upquote}}{}
\IfFileExists{microtype.sty}{% use microtype if available
  \usepackage[]{microtype}
  \UseMicrotypeSet[protrusion]{basicmath} % disable protrusion for tt fonts
}{}
\makeatletter
\@ifundefined{KOMAClassName}{% if non-KOMA class
  \IfFileExists{parskip.sty}{%
    \usepackage{parskip}
  }{% else
    \setlength{\parindent}{0pt}
    \setlength{\parskip}{6pt plus 2pt minus 1pt}}
}{% if KOMA class
  \KOMAoptions{parskip=half}}
\makeatother
\usepackage{xcolor}
\usepackage[margin=1in]{geometry}
\usepackage{graphicx}
\makeatletter
\def\maxwidth{\ifdim\Gin@nat@width>\linewidth\linewidth\else\Gin@nat@width\fi}
\def\maxheight{\ifdim\Gin@nat@height>\textheight\textheight\else\Gin@nat@height\fi}
\makeatother
% Scale images if necessary, so that they will not overflow the page
% margins by default, and it is still possible to overwrite the defaults
% using explicit options in \includegraphics[width, height, ...]{}
\setkeys{Gin}{width=\maxwidth,height=\maxheight,keepaspectratio}
% Set default figure placement to htbp
\makeatletter
\def\fps@figure{htbp}
\makeatother
\setlength{\emergencystretch}{3em} % prevent overfull lines
\providecommand{\tightlist}{%
  \setlength{\itemsep}{0pt}\setlength{\parskip}{0pt}}
\setcounter{secnumdepth}{-\maxdimen} % remove section numbering
\ifLuaTeX
  \usepackage{selnolig}  % disable illegal ligatures
\fi
\usepackage{bookmark}
\IfFileExists{xurl.sty}{\usepackage{xurl}}{} % add URL line breaks if available
\urlstyle{same}
\hypersetup{
  pdftitle={Bioaerosol Infection Risk Assessment Report},
  hidelinks,
  pdfcreator={LaTeX via pandoc}}

\title{Bioaerosol Infection Risk Assessment Report}
\usepackage{etoolbox}
\makeatletter
\providecommand{\subtitle}[1]{% add subtitle to \maketitle
  \apptocmd{\@title}{\par {\large #1 \par}}{}{}
}
\makeatother
\subtitle{Automatically generated by the \texttt{airqmra} R package}
\author{}
\date{\vspace{-2.5em}2026-05-23}

\begin{document}
\maketitle

\section{1. Executive Summary}\label{executive-summary}

This report is prepared in accordance with Chapter 6 ``Infection Risk
Assessment of Bioaerosols in the Atmospheric Environment'' of the
Technical Guideline for Health Risk Assessment of Bioaerosols in the
Atmospheric Environment (T/CSES XXXX-XXXX). It quantitatively evaluates
the infection risk posed by specific infectious pathogens under defined
exposure scenarios.

\begin{itemize}
\tightlist
\item
  \textbf{Target pathogen}: legionella pneumophila
\item
  \textbf{Exposed population}: Adult
\item
  \textbf{Microenvironment}: Outdoor, \textbf{Season}: Annual
\item
  \textbf{Simulation method}: One‑dimensional Monte Carlo (variability
  only)
\item
  \textbf{Assessment tool}: \texttt{airqmra} R package (version
  0.0.0.9000)
\end{itemize}

\section{2. Data Sources}\label{data-sources}

\begin{itemize}
\tightlist
\item
  \textbf{Exposure concentration}: Obtained from field sampling and
  laboratory analysis.
\item
  \textbf{Dose‑response parameters}: Sourced from the built‑in pathogen
  database (\texttt{pathogen\_db}) of \texttt{airqmra}, with original
  references provided in Appendix A.
\item
  \textbf{Exposure factors}: Primarily derived from the \emph{Chinese
  Exposure Factors Handbook}.
\end{itemize}

\section{3. Assessment Methods}\label{assessment-methods}

\subsection{3.1 Hazard Identification}\label{hazard-identification}

legionella pneumophila has been identified as a typical infectious
pathogen. The primary route of exposure is inhalation, which can lead to
infection in exposed populations.

\subsection{3.2 Dose‑Response Assessment}\label{doseresponse-assessment}

The dose‑response relationship is quantified using the
\textbf{exponential} model. \[P = 1 - \exp(-k \cdot d)\] Parameter
\textbf{k} = 4.99e-05

\subsection{3.3 Exposure Assessment}\label{exposure-assessment}

The inhaled dose is calculated as:
\[d = C \times IR \times T \times f_{\text{survival}} \times f_{\text{deposition}}\]

\begin{verbatim}
## **Inhalation rate (IR)**: Automatically loaded for the `r scen$exposure_group` group (long‑term breathing rate, m³/h), using a log‑normal distribution.
\end{verbatim}

\begin{verbatim}
## **Exposure time (T)**: Obtained from the time‑activity database based on microenvironment and season.
\end{verbatim}

\begin{verbatim}
## **Survival fraction (f_survival)**: Empirical log‑normal distribution (μ = -1.322, σ = 0.427).
\end{verbatim}

\begin{verbatim}
## **Deposition fraction (f_deposition)**: Pulmonary deposition for `r dep_pop` population, log‑normally distributed.
\end{verbatim}

\subsection{3.4 Risk Characterization}\label{risk-characterization}

\begin{itemize}
\tightlist
\item
  \textbf{Single‑exposure infection risk}: Directly computed from the
  dose‑response model.
\item
  \textbf{Annual infection risk}:
  \[P_{\text{annual}} = 1 - (1 - P_i)^n\] where \emph{n} is the number
  of exposure events per year.
\item
  \textbf{Uncertainty analysis}: Monte Carlo simulation combined with
  multiple sensitivity analysis methods.
\end{itemize}

\section{4. Results}\label{results}

\subsection{4.1 Single‑Exposure Infection
Risk}\label{singleexposure-infection-risk}

\begin{verbatim}
## **Summary statistics**
##      Min.   1st Qu.    Median      Mean   3rd Qu. 
## 6.173e-06 1.154e-04 2.304e-04 3.860e-04 4.538e-04 
##      Max. 
## 1.300e-02
\end{verbatim}

\includegraphics{D:/DeskTop/R/qmra/airqmra/airqmra_report_en_files/5f3501f65201ec51afb6ef8e8c28371526014fb0.pdf}

\subsection{4.2 Annual Infection Risk}\label{annual-infection-risk}

\begin{verbatim}
## **Summary statistics**
##      Min.   1st Qu.    Median      Mean   3rd Qu. 
## 0.0001235 0.0023047 0.0045975 0.0076412 0.0090366 
##      Max. 
## 0.2302102
\end{verbatim}

\includegraphics{D:/DeskTop/R/qmra/airqmra/airqmra_report_en_files/afdca20e49db8e35b9497a4c22534a9f1700b27a.pdf}

\section{5. Sensitivity Analysis}\label{sensitivity-analysis}

\subsection{5.1 Spearman Rank
Correlation}\label{spearman-rank-correlation}

\includegraphics{D:/DeskTop/R/qmra/airqmra/airqmra_report_en_files/1391765a11dfca70f00e9aa551d4630201662020.pdf}

\subsection{5.2 PRCC (Partial Rank Correlation
Coefficient)}\label{prcc-partial-rank-correlation-coefficient}

\begin{verbatim}
## PRCC analysis not provided.
\end{verbatim}

\subsection{5.3 Morris Elementary
Effects}\label{morris-elementary-effects}

\begin{verbatim}
## Morris analysis not provided.
\end{verbatim}

\subsection{5.4 Sobol' Variance
Decomposition}\label{sobol-variance-decomposition}

\begin{verbatim}
## Sobol' analysis not provided.
\end{verbatim}

\section{6. Uncertainty Analysis}\label{uncertainty-analysis}

\subsection{6.1 Propagation of Parameter
Uncertainty}\label{propagation-of-parameter-uncertainty}

\begin{verbatim}
## 2D simulation results not provided.
\end{verbatim}

\section{7. Conclusions and
Recommendations}\label{conclusions-and-recommendations}

For the scenario where Adult individuals are exposed to legionella
pneumophila in a Outdoor environment: - The median single‑exposure
infection risk is approximately \ensuremath{2.3\times 10^{-4}}. - The
median annual infection risk is approximately 0.0046. The annual risk
median exceeds the commonly acceptable benchmark (≤10⁻⁴), indicating a
potential health concern. Control measures are recommended.

\section{8. Appendices}\label{appendices}

\begin{itemize}
\tightlist
\item
  Appendix 1: Summary of main exposure and dose‑response parameters.
\item
  Appendix 2: Detailed sensitivity analysis charts.
\end{itemize}

\begin{center}\rule{0.5\linewidth}{0.5pt}\end{center}

\textbf{End of Report}

\end{document}
